\chapter{Входы машины: как к ней подключены глаза, уши и другие датчики}
\chaptermark{Входы машины}
\label{sec:sensors}

\section{Датчик как преобразователь}

На рис.~\ref{fig:machine} данные входили в машину слева, из блока «датчики». Теперь откроем этот блок. Для инженера каждый орган чувств --- это \emph{преобразователь} с аналоговым входным каскадом и аналого-цифровым преобразователем в конце, только цифры на выходе --- не числа, а импульсы.

Схема у всех органов чувств одна и та же (рис.~\ref{fig:transducer}). Физическая величина --- свет, давление, колебание, молекула --- действует на белок-рецептор в мембране чувствительной клетки. Этот белок работает как затвор: он сам или через короткую цепочку реакций открывает ионный канал. Возникает \emph{рецепторный ток}, а с ним --- аналоговое напряжение на мембране. Дальше напряжение проходит регулировку усиления и попадает в кодер импульсов --- уже знакомый интегрирующий нейрон~\eqref{eq:glutneuron}, который превращает уровень в частоту и моменты импульсов. Выход почти всегда один и тот же: чувствительные нейроны глаза, уха, обоняния и кожи выбрасывают глутамат, так что входы подключаются прямо к шине \nt{ГЛУТ}.

\begin{figure}[t]
\centering
\begin{tikzpicture}[>=Stealth,
  blk/.style={draw,very thick,rounded corners=2pt,align=center,font=\footnotesize,minimum height=1.0cm,text width=1.9cm,fill=blue!5}]
\node[align=center,font=\small] (P) at (0.1,0) {свет, звук,\\давление,\\молекула};
\node[blk] (R) at (3.0,0) {рецептор-\\затвор};
\node[blk] (G) at (6.5,0) {регулировка\\усиления};
\node[blk] (E) at (10.0,0) {кодер\\импульсов};
\node[align=center,font=\small] (B) at (13.4,0) {шина\\\nt{ГЛУТ}};
\draw[->,very thick] (P) -- (R);
\foreach \ic/\xx in {sun/-1.1,speaker/-0.3,press/0.5,molecule/1.3}{\pic[scale=0.85] at (\xx,1.05) {\ic};}
\draw[->,very thick] (R) -- node[above,font=\scriptsize] {ток} (G);
\draw[->,very thick] (G) -- node[above,font=\scriptsize] {уровень} (E);
\draw[->,very thick,red!70!black] (E) -- node[above,font=\scriptsize,black] {импульсы} (B);
\draw[decorate,decoration={brace,amplitude=5pt,mirror},gray!60!black] (1.8,-0.8) -- (7.7,-0.8) node[midway,below=5pt,font=\scriptsize,black] {аналоговая часть};
\draw[decorate,decoration={brace,amplitude=5pt,mirror},gray!60!black] (8.8,-0.8) -- (11.2,-0.8) node[midway,below=5pt,font=\scriptsize,black] {импульсы};
\end{tikzpicture}
\caption{Любой орган чувств как цепь преобразований. Белок-рецептор открывает канал и даёт ток; регулировка усиления приспосабливает его к обстановке; кодер~\eqref{eq:glutneuron} превращает уровень в импульсы, которые уходят на шину \nt{ГЛУТ}. В глазу и ухе первые каскады вообще не выдают импульсов: сигнал остаётся аналоговым, пока не приходится передавать его далеко.}
\label{fig:transducer}
\end{figure}

Одна деталь хорошо знакома электронщику. В глазу и в ухе самые первые клетки импульсов не выдают вовсе: они передают соседям плавное напряжение, как аналоговый каскад на плате. Импульсы появляются только там, где сигнал надо везти далеко. Аналоговый сигнал на расстоянии миллиметров затухает и зашумляется, а импульс, как мы видели в главе~\ref{sec:neurontypes}, доходит до конца провода той же высоты. Оцифровывают там, где начинается длинная линия.

\section{На пределе физики}

Датчики мозга работают у физических пределов чувствительности. Датчик света в темноте отвечает на \emph{один фотон}: поглощение одной частицы света даёт ток около пикоампера, заметный на фоне собственного шума~\cite{Baylor1979}. Датчик звука замечает смещение своего чувствительного пучка меньше чем на нанометр~\cite{Hudspeth1989}.

Когда света мало, точность ограничивает не датчик, а сам свет: фотоны приходят случайно, пуассоновским потоком. Если за время накопления $T$ в датчик попадает в среднем $N=\Phi T$ фотонов, их число дрожит на $\sqrt N$. Чтобы прибавка $\Delta N$ была заметна, она должна быть в несколько раз больше этого дрожания, и различимая доля изменения яркости равна
\begin{empheq}[box=\keybox]{equation}
\frac{\Delta I}{I} \;\approx\; \frac{k}{\sqrt{\Phi\,T}} .
\label{eq:photonnoise}
\end{empheq}
Это та же формула, что и для числа импульсов~\eqref{eq:rateerror}: дробовой шум фотонов и дробовой шум импульсов подчиняются одному закону. В темноте глаз может улучшить точность только одним способом --- дольше копить фотоны, и время накопления действительно растёт до десятых долей секунды. Поэтому в сумерках мы видим медленнее.

\section{Динамический диапазон: регулировка усиления}

Самая трудная инженерная задача у датчиков --- диапазон. Освещённость от звёздной ночи до солнечного снега меняется примерно на десять порядков, интенсивность звука от порога слышимости до болевого --- на двенадцать. А частота импульсов меняется от нуля до нескольких сотен в секунду, меньше чем на три порядка. Линейно такой вход в такой выход не уложить.

Решение то же, что у любого приёмника: \emph{автоматическая регулировка усиления}. Ответы чувствительных клеток глаза хорошо описываются формулой
\begin{empheq}[box=\keybox]{equation}
r \;=\; r_{\max}\,\frac{I}{I+\sigma} ,
\label{eq:nakarushton}
\end{empheq}
где $\sigma$ --- яркость, при которой ответ достигает половины~\cite{NakaRushton1966}. Сама по себе~\eqref{eq:nakarushton} охватывает лишь два-три порядка. Но $\sigma$ не постоянна: при долгой работе на ярком фоне она растёт вслед за средней яркостью. Кривая ответа сдвигается вдоль логарифмической оси яркости и всякий раз ставит свой крутой участок туда, где сейчас находится вход (рис.~\ref{fig:adaptation}). Это то же деление на общую активность, что и шунтирующее торможение главы~\ref{sec:gaba}~\cite{Carandini2012}, только знаменатель здесь --- средняя яркость за последние секунды.

\begin{figure}[t]
\centering
\begin{tikzpicture}[>=Stealth,x=1.25cm,y=2.8cm]
\draw[->,gray!70!black] (-2,0) -- (6.4,0) node[below left,font=\small,black] {$\log_{10} I$};
\draw[->,gray!70!black] (-2,0) -- (-2,1.15) node[left,font=\small,black] {$r/r_{\max}$};
\foreach \u in {-2,0,2,4,6}{\draw[gray!60!black] (\u,-0.02) -- (\u,0.02); \node[below,font=\scriptsize] at (\u,0) {$\u$};}
\draw[densely dashed,gray!60!black] (-2,0.5) -- (6.2,0.5);
\foreach \s/\c/\lab in {0/blue!60!black/{тёмный фон},2/green!45!black/{средний},4/red!70!black/{яркий}}{
  \pgfmathsetmacro{\lo}{max(-2,\s-3)}
  \draw[very thick,\c] (-2,0) -- (\lo,0) plot[domain=\lo:6.2,samples=80] (\x,{1/(1+10^(\s-\x))});
  \node[font=\scriptsize,\c,anchor=west] at (\s+0.15,0.42) {\lab};}
\foreach \s/\ic in {0/moon,2/cloud,4/sun}{\pic at (\s+0.6,0.64) {\ic};}
\end{tikzpicture}
\caption{Регулировка усиления датчика света по~\eqref{eq:nakarushton}. Каждая кривая охватывает два-три порядка яркости; при переходе на более яркий фон $\sigma$ растёт, и кривая сдвигается по логарифмической оси. Вместе сдвигающиеся кривые покрывают весь диапазон.}
\label{fig:adaptation}
\end{figure}

У регулировки есть цена, знакомая по главе~\ref{sec:code}: датчик сообщает не яркость, а яркость \emph{по отношению к фону}. Постоянный вход постепенно перестаёт вызывать ответ, а каждое изменение вызывает. Входы машины с самого начала говорят о переменах, а не об уровнях, и в этом та же экономия импульсов, что и в утверждении~\ref{prop:economy}~\cite{Barlow1961}. Отсюда и включающие и выключающие датчики главы~\ref{sec:glutcomputer}~\cite{Schiller1992}: одни сообщают, что стало ярче, другие --- что темнее.

Инженеры повторили этот приём в \emph{событийных камерах}. Такая камера не снимает кадры по часам: каждый её пиксель сам, без общего такта, выдаёт событие «ярче» или «темнее», когда логарифм яркости изменился на заданную долю~\cite{Lichtsteiner2008}. Это ровно двухпроводный код включения---выключения из главы~\ref{sec:glutcomputer}, воплощённый в кремнии, и он даёт камере диапазон и быстроту, недоступные обычной матрице~\cite{Gallego2022}.

\section{Глаз: сжатие до провода}

В глазу около $10^8$ датчиков света~\cite{Curcio1990}, а выходных нейронов, которые везут изображение в мозг, --- около миллиона. Прежде чем сигнал покинет глаз, его сжимают примерно в сто раз. Главные приёмы сжатия нам знакомы. Каждый выходной нейрон отвечает на разницу между яркостью в маленьком пятне и яркостью вокруг него --- это вычитание соседей через торможение, как в главе~\ref{sec:gaba}. Ровные участки картинки почти ничего не стоят, а края и перемены передаются. Разные выходные нейроны специализированы: одни сообщают о посветлении, другие о потемнении, третьи о движении в определённую сторону (рис.~\ref{fig:direction}); у мыши таких выходных каналов насчитали больше тридцати~\cite{Baden2016}.

Какова пропускная способность получившегося кабеля? У морской свинки по записям выходных нейронов измерили около $875$ тысяч бит в секунду на $\sim10^5$ таких нейронов; пересчёт на миллион выходных нейронов человека даёт порядка $10^7$ бит в секунду~\cite{Koch2006} --- столько же, сколько старый сетевой кабель на десять мегабит. При средней частоте в несколько импульсов в секунду на нейрон это несколько бит на импульс, как и получалось в главе~\ref{sec:code}.

\section{Ухо: банк фильтров}

Ухо решает другую задачу: звук надо разложить по частотам. Инженер поставил бы банк полосовых фильтров, и ухо делает именно это механически. Звуковое колебание бежит по упругой перегородке, свойства которой плавно меняются вдоль её длины, и каждый участок резонирует на свою частоту. Датчики, сидящие на участке, отвечают только на свою полосу (рис.~\ref{fig:filterbank}). Частоты разложены по месту примерно логарифмически: каждой октаве достаётся похожая доля датчиков. Номер сработавшего датчика --- это частота, код местом из главы~\ref{sec:code}.

\begin{figure}[t]
\centering
\begin{tikzpicture}[>=Stealth,x=1.35cm,y=2.2cm]
\draw[->,gray!70!black] (-0.3,0) -- (7.6,0) node[above left,font=\small,black] {частота, Гц};
\draw[->,gray!70!black] (-0.3,0) -- (-0.3,1.15) node[left,font=\small,black] {ответ};
\foreach \k/\f in {0/125,1/250,2/500,3/1000,4/2000,5/4000,6/8000}{
  \draw[gray!60!black] (\k,-0.02) -- (\k,0.02); \node[below,font=\scriptsize] at (\k,0) {\f};
  \draw[thick,blue!60!black] plot[domain={max(\k-1.2,-0.3)}:{\k+0.7},samples=60] (\x,{1/(1+((\x-\k)/(\x<\k?0.45:0.2))^4)});}
% a piano keyboard: seven white keys per octave, like the octaves on the axis
\foreach \i in {0,...,48}{\draw[gray!50!black,fill=white,line width=0.3pt] ({\i/7},-0.44) rectangle ({(\i+1)/7},-0.26);}
\foreach \o in {0,...,6}{\foreach \j in {0,1,3,4,5}{\fill[black] ({\o+(\j+1)/7-0.035},-0.35) rectangle ({\o+(\j+1)/7+0.035},-0.26);}}
\end{tikzpicture}
\caption{Ухо как банк полосовых фильтров, схема. Каждая кривая --- ответ датчиков одного участка на звук разной частоты; центральные частоты идут через октаву, как клавиши на логарифмической шкале. У настоящих фильтров крутой высокочастотный склон и пологий низкочастотный. Внизу --- клавиатура: на ней, как и в ухе, каждой октаве достаётся одинаковое место.}
\label{fig:filterbank}
\end{figure}

Резонаторы уха не пассивны. Часть датчиков сама раскачивает перегородку в такт звуку и усиливает слабые колебания в тысячи раз по амплитуде (на $66$--$76$\,дБ), а с ростом громкости усиление падает~\cite{Ruggero1997,Robles2001}. Для инженера это усилитель с положительной обратной связью, поставленный у самой границы самовозбуждения, --- та же жизнь на грани, что и у сети главы~\ref{sec:gaba}: у границы система особенно чувствительна к малым сигналам. Сжатие громкости здесь делается тем же усилителем: тихое усиливается сильно, громкое --- слабо.

У уха есть и второй код. На низких частотах импульсы нейронов идут в фазе со звуком: нейрон срабатывает на определённом участке каждой волны, хотя и не на каждой волне. У морской свинки привязка к фазе начинает слабеть около $600$\,Гц и ещё заметна примерно до $3.5$\,кГц~\cite{PalmerRussell1986}. Так в моментах импульсов сохраняется точное время звука, а из него --- разница прихода звука в два уха, по которой сеть главы~\ref{sec:glutcomputer} находит направление~\cite{Grothe2010}. На высоких частотах импульсы за волной не успевают, и остаётся только код местом. Одно ухо пользуется обоими кодами главы~\ref{sec:code}: временем там, где нейроны успевают, и местом там, где нет.

\section{Остальные датчики}

\begin{table}[t]
\centering
\footnotesize
\setlength{\tabcolsep}{3pt}
\renewcommand{\arraystretch}{1.15}
\begin{tabular}{>{\raggedright}p{1.9cm}>{\raggedright}p{2.6cm}>{\raggedright}p{4.0cm}>{\raggedright\arraybackslash}p{4.6cm}}
\toprule
Чувство & Что измеряет & Как устроен вход & Приём из книги \\
\midrule
зрение & свет & $\sim10^8$ датчиков, сжатие в $\sim100$ раз, $\sim10^7$ бит/с & регулировка усиления, вычитание соседей, два провода, направление движения \\
слух & давление воздуха & механический банк фильтров, активный усилитель & код местом и код временем, задержки \\
осязание & давление, вибрация & быстро и медленно привыкающие датчики & пара «фильтр верхних частот --- фильтр нижних» \\
температура & тепло & отдельные датчики тепла и холода & двухпроводный код \\
обоняние & молекулы & $\sim400$ типов рецепторов, у каждого датчика один тип & комбинаторный код: запах --- набор сработавших типов \\
вкус & вещества в пище & пять качеств: сладкое, горькое, солёное, кислое, умами & меченые линии; часть клеток выбрасывает АТФ или серотонин, а не глутамат \\
положение тела & растяжение мышц & датчики длины и скорости растяжения & обратная связь для регулятора движения \\
\bottomrule
\end{tabular}
\caption{Как подключены датчики. Все устройства в третьем столбце измерены; правый столбец связывает их с приёмами из предыдущих глав.}
\label{tab:sensors}
\end{table}

Остальные чувства собраны в таблице~\ref{tab:sensors}, и почти в каждой строке узнаётся приём из предыдущих глав. Осязание пользуется парой фильтров: одни датчики давления отвечают на прикосновение и сразу замолкают, другие держат ответ, пока давление есть. Температуру передают два провода, отдельные для тепла и для холода. Самый необычный вход --- обоняние. Типов обонятельных рецепторов у человека около четырёхсот, и каждый обонятельный нейрон несёт только один тип~\cite{Mombaerts2004}. Запах активирует не один тип, а набор, и сообщение --- это \emph{какие} типы сработали. Для инженера это двоичный вектор длиной около четырёхсот, у которого число возможных значений астрономическое.

\begin{keyblock}
\begin{proposition}[Входы машины]
\label{prop:sensors}
Каждое чувство --- преобразователь, который работает у физического предела чувствительности, сжимает огромный динамический диапазон автоматической регулировкой усиления и передаёт по шине \nt{ГЛУТ} прежде всего \emph{перемены}. Для этого датчики пользуются теми же приёмами, что и сама машина: двухпроводным кодом, кодом местом, кодом временем и делением на фон. Устройство датчиков измерено; то, что их коды подстроены под наибольшую информацию на импульс, --- хорошо обоснованная гипотеза.
\end{proposition}
\end{keyblock}

Входы, как и всё остальное, работают асинхронно. Каждый датчик выдаёт импульс, когда у него есть что сообщить, а разные чувства приходят с разными задержками: звук --- через миллисекунды, свет --- через десятки миллисекунд. Как машина сводит их в одну картину мира, если общего такта нет, --- одна из задач согласования во времени из главы~\ref{sec:synthesis}.

\section{Задачи}

\begin{exercise}[\lvlA]
\label{ex:11:1}
\emph{Сколько копить фотоны.} В сумерках в датчик света попадает $\Phi=100$ фотонов в секунду. Прибавка считается заметной, если она в $k=3$ раза больше дрожания, как в~\eqref{eq:photonnoise}.
\par\noindent(а)~Сколько времени нужно одному датчику, чтобы заметить изменение яркости на $10\,\%$?
\par\noindent(б)~Время накопления ограничено $0.2\,\text{с}$. Какое наименьшее изменение яркости заметит один датчик?
\par\noindent(в)~Сколько независимых датчиков нужно сложить, чтобы заметить $10\,\%$ за $0.2\,\text{с}$? Во сколько раз при этом падает разрешение по каждой оси, если датчики складываются квадратным пятном? Какой приём главы~\ref{sec:code} здесь повторяется?
\par\noindent(г)~Днём $\Phi=10^4$ в секунду. Сколько времени нужно одному датчику для тех же $10\,\%$?
\end{exercise}

\begin{exercise}[\lvlA]
\label{ex:11:2}
\emph{Сколько бит в импульсе.} (а)~Глаз передаёт $\sim10^7$ бит/с по $10^6$ выходным нейронам со средней частотой $3$--$5$ импульсов в секунду. Какую долю предельной ёмкости~\eqref{eq:spikeentropy} при $\Delta t=1\ms$ он использует?
\par\noindent(б)~Сколько бит в секунду остаётся на один датчик света после сжатия?
\par\noindent(в)~Обонятельный код: запах включает $20$ из $\approx400$ типов рецепторов. Сколько бит несёт такой набор, если все наборы равновероятны? Сколько общих типов в среднем у двух случайных запахов и какова вероятность, что общих типов три и больше? пять и больше?
\end{exercise}

\begin{exercise}[\lvlB]
\label{ex:11:3}
\emph{Что умеет одна кривая~\eqref{eq:nakarushton}.} (а)~В каком диапазоне яркостей ответ растёт от $10\,\%$ до $90\,\%$ от $r_{\max}$? Сколько это порядков?
\par\noindent(б)~Найдите наибольшую крутизну $\dd r/\dd\log_{10}I$ и яркость, при которой она достигается.
\par\noindent(в)~Сколько положений сдвигающейся кривой нужно, чтобы покрыть десять порядков яркости? двенадцать порядков громкости?
\par\noindent(г)~Пусть регулировка усиления довела $\sigma$ до яркости фона $I_b$. Докажите, что ответ на ступеньку $I_b\to I_b+\Delta I$ зависит только от контраста $\Delta I/I_b$ (закон Вебера). Пусть выход --- пуассоновский счёт импульсов за $T$, и ступенька заметна, когда прибавка счёта равна его разбросу. Найдите наименьший заметный контраст для одного нейрона при $r_{\max}=200$ в секунду и $T=0.1\,\text{с}$. Сколько независимых нейронов нужно для $2\,\%$? Каков предел при корреляции $c=0.1$~\eqref{eq:correlated}, если корреляция похожа на сигнал?
\end{exercise}

\begin{exercise}[\lvlB]
\label{ex:11:4}
\emph{Предел кода временем в ухе.} Разница прихода звука в два уха у человека --- до $0.22\,\text{м}/343\,\text{м/с}$ (глава~\ref{sec:glutcomputer}).
\par\noindent(а)~Импульсы идут в фазе с тоном частоты $f$, поэтому по ним разница прихода известна лишь с точностью до периода $1/f$. Докажите, что направление определяется однозначно, только если $f<1/(2\,\delta t_{\max})$, и вычислите эту частоту. Сравните с частотой, до которой импульсы идут в фазе.
\par\noindent(б)~Сколько детекторов в ряду~\eqref{eq:itd} сработает на тон $500$~Гц и на тон $2$~кГц от одного источника?
\par\noindent(в)~Каждый импульс дрожит на $0.5\ms$, а различать нужно $20$~мкс. Сколько независимых импульсов надо усреднить? Сколько нейронов, если каждый за $50\ms$ выдаёт импульсы с частотой $100$ в секунду?
\end{exercise}

\begin{exercise}[\lvlB]
\label{ex:11:5}
\emph{Проектирование: событийная камера на кабель глаза.} Нужна событийная камера из $10^6$ пикселей с выходом $10^7$ бит/с, как у глаза. Пиксель выдаёт событие «ярче» или «темнее», когда $\ln I$ изменился на $C$ с прошлого события. Событие несёт адрес пикселя и знак; время события --- момент его прихода, отдельно оно не передаётся.
\par\noindent(а)~Сколько бит в событии и сколько событий в секунду пропускает выход?
\par\noindent(б)~Камера должна замечать перепады яркости в $20\,\%$. Выберите $C$. Сколько событий даёт пиксель, через который проходит край с отношением яркостей $4$?
\par\noindent(в)~Край длиной $500$ пикселей движется по кадру. С какой наибольшей скоростью он может двигаться без переполнения выхода? Можно ли, повышая $C$, пропустить край, бегущий $1000$ пикселей в секунду?
\par\noindent(г)~Сравните с обычной камерой, которая через тот же выход передаёт полные кадры по $8$ бит на пиксель: сколько кадров в секунду и на сколько пикселей сдвинется край из (в) между кадрами? Какой приём главы~\ref{sec:glutcomputer} даёт событийной камере преимущество?
\end{exercise}

\begin{exercise}[\lvlB]
\label{ex:11:6}
\emph{Моделирование: регулировка усиления.} Для этой главы файла в \texttt{models/} нет; напишите модель на Python. Датчик отвечает по~\eqref{eq:nakarushton} с $r_{\max}=1$, а $\sigma$ подстраивается под яркость с $\tau=1\,\text{с}$ по одному из двух законов: в логарифмах, $\tau\,\dd\ln\sigma/\dd t=\ln I-\ln\sigma$, или линейно, $\tau\,\dd\sigma/\dd t=I-\sigma$. Фон $I_b=1$ в течение $10\,\text{с}$, затем $100$ в течение $10\,\text{с}$, затем снова $1$. Каждые $100\ms$ подаётся проба: яркость на $10\,\ms$ повышается на $10\,\%$. Ответ на пробу считайте как разность с таким же датчиком без проб.
\par\noindent(а)~Постройте ответы на пробы во времени для обоих законов. Каков ответ после полного привыкания? Сравните с $1.1/2.1-1/2$.
\par\noindent(б)~Каков ответ на первую пробу после скачка фона?
\par\noindent(в)~Через сколько секунд ответ на пробу восстанавливается до половины привыкшего после скачка вверх и после скачка вниз? Выведите эти времена из~\eqref{eq:nakarushton} для обоих законов и сравните.
\end{exercise}

\begin{exercise}[\lvlC]
\label{ex:11:7}
\emph{Под какой мир подстроена кривая.} Утверждение~\ref{prop:sensors} называет гипотезой то, что коды датчиков подстроены под наибольшую информацию.
\par\noindent(а)~Пусть к выходу $r(I)$ добавляется малый шум постоянной величины. Докажите, что информация о яркости наибольшая, когда распределение выхода равномерно, то есть $r(I)=r_{\max}\,\Phi_I(I)$, где $\Phi_I$ --- функция распределения яркостей.
\par\noindent(б)~Для какого распределения яркостей кривая~\eqref{eq:nakarushton} оптимальна в смысле (а)? Какой у него разброс $\log_{10}I$?
\par\noindent(в)~Если шум выхода пуассоновский, его дисперсия пропорциональна $r$. Какая кривая оптимальна тогда?
\par\noindent(г)~Открытый вопрос. Регулировка усиления в главе двигает только $\sigma$, то есть середину кривой. Предложите правило, которое подстраивало бы ещё и крутизну под разброс яркостей сцены, запишите его в непрерывном времени и проверьте на модели задачи~\ref{ex:11:6} со сценами разного контраста. Как опытом отличить датчик, который максимизирует информацию на импульс, от датчика, максимизирующего информацию в секунду?
\end{exercise}
